% !TeX root = ../数学分析教材（伍胜健）笔记.tex
% 本文件由主文档通过 \input 载入；请编译主文档。
\setcounter{chapter}{0} % 保留原书章号。
\chapter{函数}
\label{chap:functions}

\section{实数}
\label{sec:real-numbers}

\subsection{数集}
\label{subsec:number-sets}

\subsection{实数系的连续性}
\label{subsec:continuity-of-real-numbers}

\subsection{有界集与确界}
\label{subsec:bounded-sets-and-bounds}

以下均设 $E$ 为非空实数集，即 $\varnothing\neq E\subseteq\mathbb{R}$，所讨论的界均为实数。

\begin{definition}
\label{def:upper-and-lower-bounds}
若存在 $M\in\mathbb{R}$，使得
\[
    \forall x\in E,\quad x\leq M,
\]
则称 $E$ \emph{有上界}（bounded above），并称 $M$ 为 $E$ 的一个\emph{上界}（upper bound）。若存在 $m\in\mathbb{R}$，使得
\[
    \forall x\in E,\quad m\leq x,
\]
则称 $E$ \emph{有下界}（bounded below），并称 $m$ 为 $E$ 的一个\emph{下界}（lower bound）。
\end{definition}

\emph[caution]{上界和下界都不要求属于 $E$，也通常不唯一}：若 $M$ 是上界，则任意 $M'\geq M$ 也是上界；若 $m$ 是下界，则任意 $m'\leq m$ 也是下界。

\begin{definition}
\label{def:bounded-and-unbounded-sets}
若 $E$ 既有上界又有下界，则称 $E$ 为\emph{有界集}（bounded set），即存在实数 $m,M$，使得
\[
    \forall x\in E,\quad m\leq x\leq M.
\]
若 $E$ 不是有界集，则称 $E$ 为\emph{无界集}（unbounded set）。具体地，$E$ 无上界（unbounded above）是指
\[
    \forall M\in\mathbb{R},\quad \exists x\in E,\quad x>M;
\]
$E$ 无下界（unbounded below）是指
\[
    \forall m\in\mathbb{R},\quad \exists x\in E,\quad x<m.
\]
\emph[caution]{因此，$E$ 无界意味着它至少在一个方向上无界，并不要求它同时无上界和无下界。}
\end{definition}

\begin{definition}
\label{def:maximum-and-minimum}
若 $a\in E$ 且对所有 $x\in E$ 都有 $x\leq a$，则称 $a$ 为 $E$ 的\emph{最大值}（maximum），记为 $a=\max E$。若 $b\in E$ 且对所有 $x\in E$ 都有 $b\leq x$，则称 $b$ 为 $E$ 的\emph{最小值}（minimum），记为 $b=\min E$。
\end{definition}

\begin{definition}
\label{def:supremum}
若 $\beta\in\mathbb{R}$ 是 $E$ 的上界，且对 $E$ 的任意上界 $M$ 都有 $\beta\leq M$，则称 $\beta$ 为 $E$ 的\emph{上确界}（supremum），记为 $\beta=\sup E$。也就是说，上确界是\emph{最小的上界}（least upper bound）。等价地，$\beta=\sup E$ 当且仅当
\[
    \begin{cases}
        \forall x\in E,\quad x\leq\beta,\\
        \forall\varepsilon>0,\quad \exists x_{\varepsilon}\in E,\quad
        \beta-\varepsilon<x_{\varepsilon}\leq\beta.
    \end{cases}
\]
第一条说明 $\beta$ 是上界，\emph[tip]{第二条说明无论把 $\beta$ 减小多少，所得的数都不再是上界}。事实上，任意小于 $\beta$ 的实数都可写成 $\beta-\varepsilon$，其中 $\varepsilon>0$，故第二条恰好排除了所有更小的上界。
\end{definition}

\begin{definition}
\label{def:infimum}
若 $\alpha\in\mathbb{R}$ 是 $E$ 的下界，且对 $E$ 的任意下界 $m$ 都有 $m\leq\alpha$，则称 $\alpha$ 为 $E$ 的\emph{下确界}（infimum），记为 $\alpha=\inf E$。也就是说，下确界是\emph{最大的下界}（greatest lower bound）。等价地，$\alpha=\inf E$ 当且仅当
\[
    \begin{cases}
        \forall x\in E,\quad \alpha\leq x,\\
        \forall\varepsilon>0,\quad \exists y_{\varepsilon}\in E,\quad
        \alpha\leq y_{\varepsilon}<\alpha+\varepsilon.
    \end{cases}
\]
第一条说明 $\alpha$ 是下界，第二条说明无论把 $\alpha$ 增大多少，所得的数都不再是下界；因为任意大于 $\alpha$ 的实数都可写成 $\alpha+\varepsilon$，其中 $\varepsilon>0$。
\end{definition}

确界若存在，则必唯一。例如，若 $\beta_1,\beta_2$ 都是 $E$ 的上确界，由最小上界的定义，分别有 $\beta_1\leq\beta_2$ 和 $\beta_2\leq\beta_1$，故 $\beta_1=\beta_2$；下确界的唯一性同理。\emph[caution]{确界不一定属于原集合，而最大值和最小值必须属于原集合。}\emph{当上确界属于 $E$ 时，它就是最大值；当下确界属于 $E$ 时，它就是最小值}。反过来，最大值必为上确界，最小值必为下确界：例如，若 $a=\max E$，则 $a$ 是上界，且由 $a\in E$ 可知每个上界都不小于 $a$，所以 $a=\sup E$。

例如，开区间 $(0,1)$ 和闭区间 $[0,1]$ 的上确界都是 $1$，下确界都是 $0$。对任意 $\varepsilon>0$，取 $\delta=\min\{\varepsilon/2,1/2\}$，则 $\delta$ 与 $1-\delta$ 都属于 $(0,1)$，且 $0<\delta<\varepsilon$、$1-\varepsilon<1-\delta<1$，因而两个区间都满足上述确界刻画。对于 $(0,1)$，两个确界均不属于集合，所以它既无最大值也无最小值；对于 $[0,1]$，则有 $\max[0,1]=1$、$\min[0,1]=0$。

\subsection{几个常用不等式}
\label{subsec:common-inequalities}

\begin{definition}
\label{def:real-absolute-value}
实数 $x$ 的\emph{绝对值}（absolute value）定义为
\[
    |x|=
    \begin{cases}
        x, & x\geq 0,\\
        -x, & x<0,
    \end{cases}
\]
它表示数轴上 $x$ 到原点的距离，也可写成 $|x|=\sqrt{x^2}$。
\end{definition}

由定义可知，$-|x|\leq x\leq |x|$；对任意 $c\geq 0$，$|x|\leq c$ 等价于 $-c\leq x\leq c$，也等价于 $x\leq c$ 与 $-x\leq c$ 同时成立。这些性质使我们能够用绝对值统一估计正、负两种情形下的大小。

\begin{proposition}[三角不等式]
\label{prop:real-triangle-inequality}
对任意实数 $a,b$，都有
\begin{equation}
    |a+b|\leq |a|+|b|.
    \label{eq:real-triangle-inequality}
\end{equation}
式~\eqref{eq:real-triangle-inequality} 称为\emph{三角不等式}（triangle inequality），即两数之和的绝对值不超过它们绝对值之和；等号成立当且仅当 $ab\geq 0$。
\end{proposition}

\begin{proof}
由绝对值的性质，$-|a|\leq a\leq |a|$ 且 $-|b|\leq b\leq |b|$，两式相加，得到
\[
    -(|a|+|b|)\leq a+b\leq |a|+|b|.
\]
由于 $|a|+|b|\geq 0$，这等价于 $|a+b|\leq |a|+|b|$。

等号成立当且仅当 $a+b=|a|+|b|$ 或 $a+b=-(|a|+|b|)$。前一种情形等价于 $(|a|-a)+(|b|-b)=0$，由于两项均非负，又等价于 $a\geq 0$ 且 $b\geq 0$；后一种情形同理等价于 $a\leq 0$ 且 $b\leq 0$。因此，等号成立当且仅当 $ab\geq 0$，也就是 $a,b$ 同号或至少有一个为零。
\end{proof}

\emph[caution]{三角不等式给出 $|a-b|\leq |a|+|b|$ 这一上界；利用它还可得到 $|a-b|$ 的下界，即下面的反向三角不等式。}

\begin{corollary}[反向三角不等式]
\label{cor:real-reverse-triangle-inequality}
对任意实数 $a,b$，都有\emph{反向三角不等式}（reverse triangle inequality）
\begin{equation}
    \bigl||a|-|b|\bigr|\leq |a-b|.
    \label{eq:real-reverse-triangle-inequality}
\end{equation}
即两数绝对值之差的绝对值不超过它们之差的绝对值。
\end{corollary}

\begin{proof}
将 $a$ 写成 $(a-b)+b$，将 $b$ 写成 $(b-a)+a$，分别应用命题~\ref{prop:real-triangle-inequality}，得到
\[
    \begin{aligned}
        |a|&\leq |a-b|+|b|,\\
        |b|&\leq |b-a|+|a|=|a-b|+|a|.
    \end{aligned}
\]
移项后分别有 $|a|-|b|\leq |a-b|$ 和 $|b|-|a|\leq |a-b|$，合并即得
\[
    -|a-b|\leq |a|-|b|\leq |a-b|.
\]
由于 $|a-b|\geq 0$，根据绝对值的性质，这等价于 $\bigl||a|-|b|\bigr|\leq |a-b|$。
\end{proof}

\begin{definition}
\label{def:arithmetic-and-geometric-means}
设 $n$ 为正整数，$a_1,\ldots,a_n$ 为非负实数。它们的\emph{算术平均}（arithmetic mean）与\emph{几何平均}（geometric mean）分别定义为
\[
    A_n=\frac{a_1+\cdots+a_n}{n},
    \qquad
    G_n=\sqrt[n]{a_1a_2\cdots a_n}.
\]
其中，几何平均取非负的 $n$ 次方根。下面的不等式说明，这两个平均值之间存在确定的大小关系。
\end{definition}

\begin{proposition}[算术平均--几何平均不等式]
\label{prop:am-gm-inequality}
对任意正整数 $n$ 及非负实数 $a_1,\ldots,a_n$，都有
\begin{equation}
    \sqrt[n]{a_1a_2\cdots a_n}
    \leq \frac{a_1+\cdots+a_n}{n}.
    \label{eq:am-gm-inequality}
\end{equation}
式~\eqref{eq:am-gm-inequality} 称为\emph{算术平均--几何平均不等式}（arithmetic mean--geometric mean inequality，简称 AM--GM 不等式），即 $G_n\leq A_n$；\emph[caution]{等号成立当且仅当 $a_1=\cdots=a_n$。}
\end{proposition}

\begin{proof}
对 $n$ 作数学归纳。$n=1$ 时两边相等；$n=2$ 时，由
\[
    \frac{a_1+a_2}{2}-\sqrt{a_1a_2}
    =\frac{(\sqrt{a_1}-\sqrt{a_2})^2}{2}\geq 0
\]
可得结论，且等号成立当且仅当 $a_1=a_2$。

设 $n\geq 3$，并假设结论对 $n-1$ 个非负实数成立。若 $a_1,\ldots,a_n$ 中有零，则几何平均为零，不等式成立；此时等号成立当且仅当算术平均也为零，即所有 $a_i$ 均为零。以下设各 $a_i>0$，记
\[
    A=\frac{a_1+\cdots+a_n}{n}>0.
\]
若所有 $a_i$ 相等，则等号成立。若它们不全相等，由 $\sum_{i=1}^{n}a_i=nA$ 可知，必有一项小于 $A$，另一项大于 $A$；否则各项都在 $A$ 的同一侧且至少一项严格不等，其和便不等于 $nA$。重新编号后可设 $a_1<A<a_2$。

\emph[tip]{将 $a_1,a_2$ 替换为 $A$ 和 $b=a_1+a_2-A$，保持两项之和不变，同时增大两项之积。}事实上，$b=a_1+(a_2-A)>0$，且
\[
    Ab-a_1a_2=(A-a_1)(a_2-A)>0.
\]
其余 $n-1$ 个数 $b,a_3,\ldots,a_n$ 的和为 $(n-1)A$，由归纳假设，
\[
    ba_3\cdots a_n\leq A^{n-1}.
\]
由于各项均为正，因而
\[
    a_1a_2\cdots a_n
    <Ab\,a_3\cdots a_n
    \leq A\cdot A^{n-1}=A^n.
\]
两边取非负的 $n$ 次方根，便得 $G_n<A=A_n$。因此，各数不全相等时为严格不等式，全相等时为等式，归纳完成。
\end{proof}

下面讨论幂与一次式之间的比较。先回顾两个代数恒等式：设 $a,b\in\mathbb{R}$，$m$ 为正整数，则有\emph{幂差公式}
\begin{equation}
    a^m-b^m=(a-b)\sum_{k=0}^{m-1}a^{m-1-k}b^k;
    \label{eq:integer-power-difference}
\end{equation}
当 $m$ 为奇数时，还有\emph{幂和公式}
\begin{equation}
    a^m+b^m=(a+b)\sum_{k=0}^{m-1}(-1)^k a^{m-1-k}b^k.
    \label{eq:odd-integer-power-sum}
\end{equation}
在幂差公式中，将右端乘开后，中间项 $a^{m-k}b^k$（$1\leq k\leq m-1$）分别出现一次正号和一次负号，抵消后只剩 $a^m-b^m$；幂和公式则由幂差公式中以 $-b$ 代替 $b$ 得到，其中用到 $m$ 为奇数时 $(-b)^m=-b^m$。下面利用幂差公式证明伯努利不等式。

\begin{proposition}[伯努利不等式]
\label{prop:bernoulli-inequality}
对任意实数 $x\geq-1$ 和任意正整数 $n$，都有
\begin{equation}
    (1+x)^n\geq1+nx.
    \label{eq:bernoulli-inequality}
\end{equation}
式~\eqref{eq:bernoulli-inequality} 称为\emph{伯努利不等式}（Bernoulli's inequality）；等号成立当且仅当 $n=1$ 或 $x=0$。
\end{proposition}

\begin{proof}
当 $n=1$ 或 $x=0$ 时，两边相等。以下设 $n\geq2$ 且 $x\neq0$。若 $x=-1$，则 $(1+x)^n=0>1-n=1+nx$。以下只需讨论 $x>-1$ 的情形，由幂差公式~\eqref{eq:integer-power-difference}，
\[
    (1+x)^n-1
    =x\bigl[1+(1+x)+(1+x)^2+\cdots+(1+x)^{n-1}\bigr].
\]
\emph[tip]{先比较方括号内的 $n$ 项之和与 $n$ 的大小，再根据 $x$ 的正负确定乘积的不等号方向。}当 $x>0$ 时，$1+x>1$，方括号内除首项等于 $1$ 外，其余 $n-1$ 项均大于 $1$，故其和大于 $n$；乘以 $x>0$，得到 $(1+x)^n-1>nx$。当 $-1<x<0$ 时，$0<1+x<1$，方括号内除首项外，其余各项均小于 $1$，故其和小于 $n$；由于 $x<0$，乘以 $x$ 后不等号反向，仍得到 $(1+x)^n-1>nx$。因此不等式在全部条件下成立，且 $n\geq2$、$x\neq0$ 时均为严格不等式，这也给出了所述等号条件。
\end{proof}

为了便于应用，还可通过代换 $x=B/A$，将伯努利不等式写成双参数形式。

\begin{corollary}[伯努利不等式的双参数形式]
\label{cor:two-parameter-bernoulli}
设 $A,B\in\mathbb{R}$，满足 $A>0$ 且 $A+B\geq0$。对任意正整数 $n$，都有
\begin{equation}
    (A+B)^n\geq A^n+nA^{n-1}B.
    \label{eq:two-parameter-bernoulli}
\end{equation}
当 $n=1$ 时恒取等号；当 $n>1$ 时，等号成立当且仅当 $B=0$。
\end{corollary}

\begin{proof}
由 $A>0$ 且 $A+B\geq0$，有 $1+B/A\geq0$，即 $B/A\geq-1$。在式~\eqref{eq:bernoulli-inequality} 中取 $x=B/A$，再将两边乘以正数 $A^n$，得到
\[
    (A+B)^n=A^n\left(1+\frac BA\right)^n
    \geq A^n\left(1+n\frac BA\right)
    =A^n+nA^{n-1}B.
\]
上述变形保持等号条件，而 $B/A=0$ 等价于 $B=0$，故等号条件随之成立。
\end{proof}

伯努利不等式把幂与一次式联系起来。下面两个例题分别展示换元与取商的用法：前者用于证明根式不等式，后者用于比较底数与指数同时变化的两个幂。

\begin{example}
\label{ex:bernoulli-root-bound}
设 $x\geq-1$，$n$ 为正整数。证明
\begin{equation}
    \sqrt[n]{1+x}\leq1+\frac xn,
    \label{eq:bernoulli-root-bound}
\end{equation}
并确定等号条件。这里根式取非负的 $n$ 次方根。
\end{example}

\begin{solve}
\emph[tip]{把根式设为 $1+t$，就能使原来的正整数次幂不等式适用。}令
\[
    t=\sqrt[n]{1+x}-1\geq-1.
\]
由伯努利不等式，
\[
    1+x=(1+t)^n\geq1+nt.
\]
移项并除以正整数 $n$，得到 $t\leq x/n$，即式~\eqref{eq:bernoulli-root-bound}。等号成立当且仅当 $n=1$ 或 $t=0$，而 $t=0$ 等价于 $x=0$。
\end{solve}

\begin{example}
\label{ex:bernoulli-adjacent-powers}
对任意正整数 $n$，证明
\[
    \left(1+\frac1{n+1}\right)^{n+1}
    >\left(1+\frac1n\right)^n.
\]
\end{example}

\begin{solve}
两边的底数与指数同时变化，\emph[tip]{取商后将相同指数的部分合并，可以得到适合伯努利不等式的形式。}具体地，
\[
    \frac{\left(1+\frac1{n+1}\right)^{n+1}}
         {\left(1+\frac1n\right)^n}
    =\frac{n+2}{n+1}
      \left(\frac{n(n+2)}{(n+1)^2}\right)^n
    =\frac{n+2}{n+1}
      \left(1-\frac1{(n+1)^2}\right)^n.
\]
由于 $-1/(n+1)^2\geq-1$，由伯努利不等式可得
\[
    \begin{aligned}
        \frac{\left(1+\frac1{n+1}\right)^{n+1}}
             {\left(1+\frac1n\right)^n}
        &\geq\frac{n+2}{n+1}
          \left(1-\frac{n}{(n+1)^2}\right)\\
        &=\frac{(n+2)(n^2+n+1)}{(n+1)^3}
          =1+\frac1{(n+1)^3}>1.
    \end{aligned}
\]
分母为正，故原不等式成立。这一比较说明，$(1+1/n)^n$ 随正整数 $n$ 严格增大。
\end{solve}

\subsection{常用记号}
\label{subsec:common-notation}

\section{函数的概念}
\label{sec:concept-of-functions}

\subsection{函数的定义}
\label{subsec:definition-of-functions}

\subsection{由已知函数构造新函数的方法}
\label{subsec:constructing-functions}

\section{函数的性质}
\label{sec:properties-of-functions}

以下讨论实值函数 $f:D\to\mathbb{R}$，其中 $\varnothing\neq D\subseteq\mathbb{R}$ 为定义域。有界性描述函数值的范围，单调性描述函数值随自变量变化的方向，周期性与奇偶性则分别描述平移和反射所保持的关系。讨论函数性质时，必须同时明确函数的对应法则与讨论范围。

\subsection{函数的有界性}
\label{subsec:boundedness-of-functions}

设 $E$ 是 $D$ 的非空子集，记 $f(E)=\{f(x):x\in E\}$。函数 $f$ 在 $E$ 上的有界性，就是函数值集合 $f(E)$ 的有界性，因而可沿用第~\ref{subsec:bounded-sets-and-bounds} 小节中的概念。

\begin{definition}
\label{def:bounded-functions}
若存在实数 $U$，使得对所有 $x\in E$ 都有 $f(x)\leq U$，则称 $f$ 在 $E$ 上\emph{有上界}（bounded above）；若存在实数 $L$，使得对所有 $x\in E$ 都有 $L\leq f(x)$，则称 $f$ 在 $E$ 上\emph{有下界}（bounded below）。若 $f$ 在 $E$ 上既有上界又有下界，则称它在 $E$ 上\emph{有界}（bounded），等价地，
\[
    \exists M\geq 0,\quad \forall x\in E,\quad |f(x)|\leq M.
\]
若不存在这样的 $M$，则称 $f$ 在 $E$ 上\emph{无界}（unbounded），即
\[
    \forall M>0,\quad \exists x\in E,\quad |f(x)|>M.
\]
当 $E=D$ 时，简称 $f$ 为有界函数或无界函数。
\end{definition}

上述等价性来自绝对值的性质：若 $L\leq f(x)\leq U$，取 $M=\max\{|L|,|U|\}$，便有 $|f(x)|\leq M$；反之，$|f(x)|\leq M$ 给出 $-M\leq f(x)\leq M$。\emph[caution]{这里的界 $M$ 必须对所有 $x\in E$ 同时有效，不能随 $x$ 改变。}无界并不意味着某一点的函数值是无穷大，而是函数值的绝对值能够超过任意预先给定的正数。

\begin{example}
\label{ex:function-boundedness-domain}
判断函数 $f(x)=1/x$ 在 $(0,1)$ 与 $[1,+\infty)$ 上的有界性。
\end{example}

\begin{solve}
在 $(0,1)$ 上，$f(x)>1$，所以 $f$ 有下界。对任意 $M>0$，取 $x=1/(M+1)\in(0,1)$，便有 $f(x)=M+1>M$，故 $f$ 无上界，从而无界。在 $[1,+\infty)$ 上，由 $0<f(x)\leq 1$ 可知 $f$ 有界。这也说明，定义域有界并不保证函数有界，定义域无界也不妨碍函数有界。
\end{solve}

有界性只要求存在统一的界，不要求函数达到这些界。例如，$f(x)=x$ 在 $(0,1)$ 上的函数值集合仍为 $(0,1)$，其上确界为 $1$、下确界为 $0$，但函数没有最大值和最小值。

\subsection{函数的单调性}
\label{subsec:monotonicity-of-functions}

\begin{definition}
\label{def:monotone-functions}
设 $\varnothing\neq E\subseteq D$。若对任意 $x_1,x_2\in E$，由 $x_1<x_2$ 总能推出 $f(x_1)\leq f(x_2)$，则称 $f$ 在 $E$ 上\emph{单调增加}或\emph{不减}（nondecreasing）；若总能推出 $f(x_1)\geq f(x_2)$，则称 $f$ 在 $E$ 上\emph{单调减少}或\emph{不增}（nonincreasing）。这两类函数统称为\emph{单调函数}（monotone function）。

若对任意 $x_1<x_2$ 都有 $f(x_1)<f(x_2)$，则称 $f$ 在 $E$ 上\emph{严格单调增加}（strictly increasing）；若都有 $f(x_1)>f(x_2)$，则称它在 $E$ 上\emph{严格单调减少}（strictly decreasing），二者统称为严格单调函数。这里各条件中的 $x_1,x_2$ 均取自 $E$。
\end{definition}

本笔记用“单调增加、单调减少”表示允许相等的情形，用“严格”排除相等。因此，常值函数既不减又不增；严格单调函数在不同自变量处的函数值必不相同。\emph[caution]{单调性要求比较讨论范围内的任意两点，不能仅比较若干点，也不能把各个区间上的单调性直接拼接为整个定义域上的单调性。}

\begin{proposition}
\label{prop:monotone-function-endpoint-bounds}
若 $a<b$ 且 $f$ 在闭区间 $[a,b]$ 上单调，则 $f$ 在该区间上有界。具体地，若 $f$ 不减，则
\[
    f(a)\leq f(x)\leq f(b),\qquad x\in[a,b];
\]
若 $f$ 不增，则
\[
    f(b)\leq f(x)\leq f(a),\qquad x\in[a,b].
\]
因此，它在该区间上的最大值和最小值都可在端点取得。
\end{proposition}

\begin{proof}
对每个 $x\in[a,b]$，都有 $a\leq x\leq b$；分别对 $a,x$ 与 $x,b$ 应用单调性的定义，即得相应不等式，端点处允许等号。因此，两个端点值分别给出函数的最大值和最小值。
\end{proof}

\begin{example}
\label{ex:monotonicity-on-subsets}
讨论 $f(x)=x^2$ 与 $g(x)=1/x$ 的单调性。
\end{example}

\begin{solve}
对 $x_1<x_2$，有
\[
    x_2^2-x_1^2=(x_2-x_1)(x_2+x_1).
\]
若 $0\leq x_1<x_2$，右端为正；若 $x_1<x_2\leq 0$，右端为负。因此，$x^2$ 在 $[0,+\infty)$ 上严格增加，在 $(-\infty,0]$ 上严格减少，在整个 $\mathbb{R}$ 上不单调。

对于 $g$，若 $x_1<x_2$ 同在 $(-\infty,0)$ 或同在 $(0,+\infty)$，则
\[
    g(x_2)-g(x_1)=\frac{x_1-x_2}{x_1x_2}<0.
\]
所以 $g$ 在这两个区间上分别严格减少。但 $-1<1$ 而 $g(-1)=-1<1=g(1)$，不符合整个定义域上单调减少的要求；结合它在 $(0,+\infty)$ 上严格减少可知，$g$ 在 $\mathbb{R}\setminus\{0\}$ 上并不单调。
\end{solve}

\subsection{函数的周期性}
\label{subsec:periodicity-of-functions}

\begin{definition}
\label{def:periodic-functions}
若存在常数 $T>0$，使得对每个 $x\in D$ 都有
\[
    x+T\in D,\qquad x-T\in D,\qquad f(x+T)=f(x),
\]
则称 $f$ 为\emph{周期函数}（periodic function），称 $T$ 为它的一个\emph{周期}（period）。本节只讨论正周期；若所有正周期中存在最小者，则称它为 $f$ 的\emph{最小正周期}或\emph{基本周期}（fundamental period）。
\end{definition}

周期性表示函数的整个图像沿水平方向平移 $T$ 后保持不变。定义中要求定义域对正、负两个方向的平移都保持不变，不能仅在 $x$ 与 $x+T$ 恰好都属于定义域时比较函数值。由于把 $x-T$ 代入周期关系可得 $f(x-T)=f(x)$，反复使用该关系便有
\[
    f(x+kT)=f(x),\qquad x\in D,\quad k\in\mathbb{Z}.
\]
特别地，$T$ 的任意正整数倍仍是周期，但周期不一定是最小正周期。例如，定义在 $\mathbb{R}$ 上的常值函数以每个正数为周期，因而没有最小正周期。

\begin{example}
\label{ex:periods-of-sine-functions}
求 $f(x)=\sin x$ 与 $g(x)=\sin^2 x$ 在 $\mathbb{R}$ 上的最小正周期。
\end{example}

\begin{solve}
由 $\sin(x+2\pi)=\sin x$ 可知 $2\pi$ 是 $f$ 的周期。若 $T>0$ 是任意周期，令 $x=0$，得 $\sin T=0$，所以 $T=k\pi$，其中 $k$ 为正整数；再令 $x=\pi/2$，得 $\cos T=1$，所以 $k$ 必为偶数。因此，$f$ 的最小正周期为 $2\pi$。

由 $\sin(x+\pi)=-\sin x$ 可知 $\sin^2(x+\pi)=\sin^2 x$，所以 $\pi$ 是 $g$ 的周期。若 $T>0$ 是 $g$ 的周期，令 $x=0$，便有 $\sin^2 T=0$，从而 $T=k\pi$，其中 $k$ 为正整数。因此，$g$ 的最小正周期为 $\pi$。
\end{solve}

\begin{proposition}
\label{prop:periodic-function-boundedness}
设 $f:\mathbb{R}\to\mathbb{R}$ 以 $T>0$ 为周期，则 $f$ 在 $\mathbb{R}$ 上有界，当且仅当它在 $[0,T)$ 上有界。
\end{proposition}

\begin{proof}
若 $f$ 在 $\mathbb{R}$ 上有界，同一个界也适用于 $[0,T)$。反之，设 $|f(t)|\leq M$ 对所有 $t\in[0,T)$ 成立。任意实数 $x$ 都可写成 $x=kT+t$，其中 $k\in\mathbb{Z}$、$0\leq t<T$；由周期性，$|f(x)|=|f(t)|\leq M$，故 $f$ 在 $\mathbb{R}$ 上有界。
\end{proof}

\subsection{函数的奇偶性}
\label{subsec:parity-of-functions}

\begin{definition}
\label{def:even-and-odd-functions}
设定义域 $D$ \emph{关于原点对称}，即 $x\in D$ 时总有 $-x\in D$。若
\[
    \forall x\in D,\quad f(-x)=f(x),
\]
则称 $f$ 为\emph{偶函数}（even function）；若
\[
    \forall x\in D,\quad f(-x)=-f(x),
\]
则称 $f$ 为\emph{奇函数}（odd function）。若上述两类条件都不满足，则称 $f$ 为非奇非偶函数；定义域不关于原点对称的函数也不属于奇函数或偶函数。
\end{definition}

偶函数的图像关于 $y$ 轴对称，奇函数的图像关于原点对称，因为图像上的点 $(x,f(x))$ 分别对应于 $(-x,f(x))$ 与 $(-x,-f(x))$。\emph[caution]{判断奇偶性应先检查定义域的对称性，再比较 $f(-x)$ 与 $f(x)$。}若奇函数在 $0$ 处有定义，则 $f(0)=-f(0)$，所以 $f(0)=0$；但奇函数的定义域不一定包含 $0$，例如 $f(x)=1/x$。在同一非空对称定义域上，一个函数既奇又偶，当且仅当它恒为零，因为这时 $f(x)=f(-x)=-f(x)$。

\begin{example}
\label{ex:parity-and-domain}
判断定义在 $\mathbb{R}$ 上的 $f(x)=x^2$、$g(x)=x^3$、$h(x)=x^2+x$ 的奇偶性，并讨论将 $f$ 限制在 $[0,+\infty)$ 上后的情形。
\end{example}

\begin{solve}
三个函数的定义域 $\mathbb{R}$ 都关于原点对称。由 $f(-x)=x^2=f(x)$ 与 $g(-x)=-x^3=-g(x)$，可知 $f$ 是偶函数，$g$ 是奇函数。对于 $h$，有 $h(-1)=0$、$h(1)=2$，既不满足 $h(-1)=h(1)$，也不满足 $h(-1)=-h(1)$，所以它非奇非偶。将 $f$ 限制在 $[0,+\infty)$ 上后，定义域不再关于原点对称，因此该限制函数非奇非偶，尽管它仍由公式 $x^2$ 给出。
\end{solve}

\begin{proposition}[奇偶分解]
\label{prop:even-odd-decomposition}
设 $D$ 非空且关于原点对称，则任意函数 $f:D\to\mathbb{R}$ 都能唯一写成一个偶函数与一个奇函数之和：
\begin{equation}
    f(x)=f_{\mathrm e}(x)+f_{\mathrm o}(x),\qquad
    f_{\mathrm e}(x)=\frac{f(x)+f(-x)}{2},\qquad
    f_{\mathrm o}(x)=\frac{f(x)-f(-x)}{2}.
    \label{eq:even-odd-decomposition}
\end{equation}
其中两个分量的定义域均为 $D$。
\end{proposition}

\begin{proof}
由 $D$ 的对称性，式~\eqref{eq:even-odd-decomposition} 中各函数均有定义。直接计算得
\[
    f_{\mathrm e}(-x)=\frac{f(-x)+f(x)}{2}=f_{\mathrm e}(x),
    \qquad
    f_{\mathrm o}(-x)=\frac{f(-x)-f(x)}{2}=-f_{\mathrm o}(x),
\]
故 $f_{\mathrm e}$ 为偶函数，$f_{\mathrm o}$ 为奇函数，且两式相加得到 $f=f_{\mathrm e}+f_{\mathrm o}$。若另有 $f=u+v$，其中 $u$ 为偶函数、$v$ 为奇函数，则
\[
    f(x)=u(x)+v(x),\qquad f(-x)=u(x)-v(x).
\]
相加与相减分别给出 $u=f_{\mathrm e}$、$v=f_{\mathrm o}$，所以分解唯一。
\end{proof}

\section{初等函数}
\label{sec:elementary-functions}

\section*{习题一}
\addcontentsline{toc}{section}{习题一}
